\documentclass[11pt]{article}
\usepackage[T1]{fontenc}
\usepackage{textcomp}
\usepackage[margin=0.75in]{geometry}
\usepackage{amsmath,amssymb,amsthm}
\usepackage{array,booktabs}
\usepackage{hyperref}
\setlength{\parindent}{0pt}
\newtheorem{theorem}{Theorem}
\theoremstyle{definition}
\newtheorem{definition}{Definition}
\title{Flow Proof Artifact: List-append-empty-right}
\author{Generated by \texttt{flow doc proof}}
\date{Flow Proof Book}
\begin{document}
\maketitle
Appending the empty list on the right changes nothing.\\[0.5em]
\textbf{Source.} church — https://en.wikipedia.org/wiki/List\_(abstract\_data\_type)\\[0.5em]
\bigskip
\noindent{\large \textbf{Derived fact 1.} \textit{xs ++ nil = xs} \hfill List \cdot append \cdot empty is the right identity}\par
\smallskip\noindent\textit{Coordinate: List \cdot append \cdot empty is the right identity}\par
\smallskip\noindent\textit{Source: church/induction}\par
\smallskip\noindent\textit{Built on: empty is the left identity, for append on List, singleton prepends the head, for append on List}\par
\medskip
\noindent\textbf{Goal.} xs ++ nil = xs\par
\noindent$\displaystyle \forall xs \in List\quad xs ++ nil = xs$\par
\medskip
\noindent
\begin{tabular}{@{} >{\raggedright\arraybackslash}p{0.06\textwidth} >{\raggedright\arraybackslash}p{0.47\textwidth} >{\raggedleft\arraybackslash}p{0.41\textwidth} @{}}
\textbf{\#} & \textbf{Proof} & \textbf{Mathematics} \\
\midrule
\textbf{1.} & We prove that empty is the right identity for append on List. & \\[0.45em]
\textbf{2.} & We proceed by induction on xs: first the base case, then the inductive step. & \\[0.45em]
\textbf{3.} & Consider the base case in which xs  equals  nil. & \\[0.45em]
\textbf{4.} & We invoke the definitional clause governing append on List: empty is the left identity, for append on List (instantiated for nil). & \\[0.45em]
\textbf{5.} & From step 3 and step 4, we can deduce that xs plus plus nil equals xs. This establishes the base case (see step 3 and step 4). Hence proven. & $\displaystyle xs ++ nil = xs$ \\[0.45em]
\textbf{6.} & For the inductive step, suppose the claim holds for all smaller values. & \\[0.45em]
\textbf{7.} & Under the supposition in step 6, let x = head(xs). & \\[0.45em]
\textbf{8.} & Under the supposition in step 6, let rest = tail(xs). & \\[0.45em]
\textbf{9.} & Under the supposition in step 6, we cross the inductive boundary: assume the claim holds for rest (the induction hypothesis). & $\displaystyle rest ++ nil = rest$ \\[0.45em]
\textbf{10.} & We invoke the definitional clause governing append on List: singleton prepends the head, for append on List (instantiated for x, nil). & \\[0.45em]
\textbf{11.} & From step 6, step 7, step 8, step 9, and step 10, this implies xs plus plus nil equals xs. Hence proven. & $\displaystyle xs ++ nil = xs$ \\[0.45em]
\bottomrule
\end{tabular}
\label{thm:List-append-empty_is_the_right_identity}
\par\medskip\hrule
\end{document}
